\chapter{Analyse et spécification des besoins}

\section{Introduction}

Ce chapitre transforme la problématique présentée au chapitre~1 en besoins et exigences pour le système Op Support IA. Le cahier des charges ne décrivait pas un logiciel interne déjà en production chez SINAPS TEC, mais un \textbf{produit à construire} : canal de chat, agent IA, workflow, multi-agents, contenus riches, clôture et satisfaction, ainsi que tableau de bord. L'existant fonctionnel au début du stage était donc \textbf{nul} : aucune base de code héritée n'a été fournie et l'application a été écrite intégralement pendant le PFA.

Les outils du marché (chats de type Slack, widgets de support et assistants RAG) constituent uniquement des \textbf{repères d'expérience utilisateur} pour le chat \og Slack light \fg{} demandé dans le sujet ; ils ne représentent pas un système existant de SINAPS TEC. L'analyse ci-dessous distingue les besoins du système, les exigences retenues et leur statut dans le dépôt actuel.

\section{Identification des acteurs}

Quatre rôles ou composants interviennent dans le fonctionnement du système. Le tableau~\ref{tab:acteurs} présente les acteurs humains ; le diagramme de cas d'utilisation du chapitre~3 reprend les acteurs humains.

\begin{table}[H]
\centering
\caption{Acteurs du système}
\label{tab:acteurs}
\begin{tabularx}{\textwidth}{@{}lX@{}}
\toprule
\textbf{Acteur} & \textbf{Rôle} \\
\midrule
Client & Utilisateur final du support. S'identifie par Google ou par nom et e-mail. Consulte et alimente \textbf{sa} conversation. \\
Agent de support & Compte créé par inscription, \textbf{validé} par l'admin. Traite les fils escaladés. \\
Administrateur & Même modèle de compte qu'un agent, avec le rôle \texttt{admin}. Valide ou refuse les agents, consulte stats et historique. \\
\bottomrule
\end{tabularx}
\end{table}

Le système IA n'est \textbf{pas} un acteur UML au sens d'un utilisateur humain : c'est un \textbf{composant interne} qui interroge la base de connaissances, peut appeler Gemini et produit des messages de type \texttt{sender = ia}. Le backend, MongoDB et Socket.io sont également des composants techniques internes, et non des acteurs humains supplémentaires.

\section{Besoins fonctionnels}

La solution retenue est une application web exposant trois espaces, complétée par un client mobile connecté au même backend :

\begin{enumerate}[leftmargin=1.6em]
  \item \textbf{Client} (\texttt{/}) : identification, conversation unique active, réponses IA, pièces jointes, escalade / retour IA, clôture et satisfaction ;
  \item \textbf{Agent} (\texttt{/agent}) : connexion e-mail / mot de passe, file des conversations, réponses humaines, clôture ;
  \item \textbf{Administrateur} (\texttt{/admin}) : validation des inscriptions agents, statistiques, historique filtrable.
\end{enumerate}

Les besoins fonctionnels couvrent donc l'authentification des utilisateurs, la création et la consultation des conversations, l'envoi et la réception des messages, les réponses IA fondées sur la base de connaissances, l'escalade vers un agent humain et le retour vers l'IA. Ils comprennent également la gestion des fils par les agents, les pièces jointes, la clôture avec satisfaction, la validation des comptes agents, l'historique, les filtres et les statistiques du tableau de bord.

Le client mobile Expo / React Native partage les routes HTTP et le serveur Socket.io du backend. Il couvre l'authentification, le chat, les réponses IA, l'escalade humaine, la clôture avec satisfaction, les pièces jointes, le temps réel, ainsi que des écrans agent et administrateur. Les différences avec le web et le statut partiel de cette réalisation sont précisés dans le tableau des exigences.

Côté serveur, une API REST Express persiste les entités dans MongoDB, appelle un service RAG + Gemini et notifie les clients via Socket.io. La conception de cette architecture est détaillée au chapitre~3 ; la présente section se limite à l'expression des besoins.

\section{Besoins non fonctionnels}

\begin{description}[style=nextline,leftmargin=1.2em,font=\bfseries]
  \item[RNF01 --- Utilisabilité] Interface en français, simple, adaptée à un usage type Slack light (exigence de design du sujet).
  \item[RNF02 --- Adaptabilité d'écran] Fonctionnement sur desktop, tablette et smartphone via une UI responsive (Next.js) ; un client mobile Expo / React Native distinct est également présent.
  \item[RNF03 --- Sécurité] Mots de passe agents hachés (bcrypt) ; JWT ; vérification du jeton Google ; restriction d'accès aux conversations (propriétaire client ou agent/admin) ; messages \texttt{ia} non injectables depuis l'extérieur ; upload authentifié et filtré.
  \item[RNF04 --- Disponibilité du bot] Si \texttt{GEMINI\_API\_KEY} est absente ou l'appel échoue, un repli sur le meilleur extrait RAG (ou un message d'invitation à l'escalade) évite un silence total.
  \item[RNF05 --- Temps réel] Diffusion des messages et des mises à jour de conversation par Socket.io (salles par conversation et événements globaux).
  \item[RNF06 --- Testabilité] Suite Jest / Supertest / MongoMemoryServer sur le backend.
  \item[RNF07 --- Portabilité de configuration] Secrets et URI via variables d'environnement ; exemples \texttt{.env.example}.
\end{description}

Les besoins non fonctionnels décrits ici restent limités aux propriétés effectivement présentes dans le sujet ou dans le dépôt ; aucune garantie de performance, de disponibilité en production ou de confidentialité supplémentaire n'est avancée sans résultat vérifié.

\section{Exigences et périmètre}

Les exigences ci-dessous sont dérivées du document de sujet \textit{Projet Stage -- Op Support IA} et de ce qui a été effectivement spécifié pour le MVP. Le tableau~\ref{tab:exigences} indique le statut \textbf{dans le dépôt actuel}, afin de ne pas confondre le souhait initial et le réalisé.

\small
\begin{longtable}{@{}p{1.2cm}p{9.4cm}p{5.3cm}@{}}
\caption{Exigences fonctionnelles et statut dans le dépôt}
\label{tab:exigences}\\
\toprule
\textbf{Id} & \textbf{Exigence} & \textbf{Statut} \\
\midrule
\endfirsthead
\multicolumn{3}{c}{\tablename\ \thetable\ -- suite} \\
\toprule
\textbf{Id} & \textbf{Exigence} & \textbf{Statut} \\
\midrule
\endhead
\midrule
\multicolumn{3}{r}{Suite à la page suivante} \\
\endfoot
\bottomrule
\endlastfoot
RF01 & Le client identifie une demande dans une interface de chat. & Réalisé \\
RF02 & Authentification client par compte Google (OAuth). & Réalisé sur le Web et côté serveur (jeton Google vérifié). Une saisie nom/e-mail existe en parallèle ; Google Sign-In natif mobile n'a pas été vérifié. \\
RF03 & L'agent IA répond à partir d'une base de connaissances (RAG). & Réalisé (retriever TF-IDF + Gemini, repli local). \\
RF04 & Workflow de statuts (en cours, en attente, résolu) et mode IA / humain. & Réalisé \\
RF05 & Bascule vers un opérateur humain à l'initiative du client. & Réalisé \\
RF06 & Retour vers l'IA (désescalade). & Réalisé \\
RF07 & Plusieurs agents (comptes distincts, compétences saisies à l'inscription). & Comptes : réalisé. Routage automatique par compétence : \textbf{non réalisé}. \\
RF08 & Chat avec une UX de type messagerie (bulles, emojis, interface soignée). & Réalisé (web), avec indicateur de saisie, verrouillage pendant la réponse IA et retours d'action. \\
RF09 & Échange de documents, images, vidéos. & \textbf{Partiellement réalisé} : upload et échange de fichiers pris en charge par le backend et le Web ; images et documents disponibles sur mobile. Le flux vidéo mobile n'a pas été implémenté ni vérifié, et l'aperçu \og unfurl \fg{} de liens n'est pas documenté comme fonctionnalité distincte. \\
RF10 & Clôture et enquête de satisfaction (note 1--5, commentaire). & Réalisé \\
RF11 & Inscription agent, validation ou rejet par l'admin. & Réalisé \\
RF12 & Historique des conversations, recherche et filtre par statut. & Réalisé (espace admin). \\
RF13 & Statistiques : volume, résolution IA vs humain, satisfaction moyenne, temps de réponse. & Réalisé (agrégations MongoDB, pas de chiffres marketing fictifs). \\
RF14 & Application mobile native (React Native / Flutter). & \textbf{Partiellement réalisé} : client Expo / React Native présent et connecté au backend ; aucune version Flutter ni recette formelle documentée. \\
RF15 & File d'attente / messagerie temps réel. & Réalisé (Socket.io). \\
RF16 & Suggestions d'actions rapides après une réponse IA. & Réalisé (confirmation, nouvelle question, escalade ; actions contextualisées). \\
RF17 & Affichage cohérent des messages et des profils. & Réalisé (horodatages persistés, avatars client de repli déterministes, avatars agent). \\
\end{longtable}
\normalsize

Les chiffres illustratifs du document de sujet (satisfaction 94~\%, 2,4~min, 68~\% de résolutions IA, note 4,7/5) sont des \textbf{exemples de communication produit}. Ils \textbf{ne sont pas} des résultats mesurés pendant le stage et ne seront pas présentés comme tels.

\section{Contraintes et limites du périmètre}

\subsection{Contraintes}

\begin{itemize}[leftmargin=1.4em]
  \item Durée limitée (environ deux mois et demi) et travail \textbf{individuel}.
  \item Stack imposée / validée : web React (réalisé avec Next.js), backend Node/Express, MongoDB, Gemini, JWT --- et non Vue, Angular, OpenAI par défaut, ni Kafka.
  \item Hébergement applicatif : \textbf{local} au moment de la rédaction ; base distante documentée avec \textbf{MongoDB Atlas}.
  \item Base de connaissances : jeu d'\textbf{exemples} (commandes, remboursements, mot de passe, codes promo), pas un corpus métier SINAPS TEC fourni par ailleurs.
\end{itemize}

\subsection{Hors périmètre (à ce stade)}

\begin{itemize}[leftmargin=1.4em]
  \item application iOS / Android native ;
  \item file de messages (Kafka, RabbitMQ) ;
  \item RAG vectoriel (embeddings, base vectorielle, LangChain) ;
  \item affectation automatique d'un agent selon ses compétences ;
  \item configuration complète du déploiement public non versionnée dans le dépôt : le Web public et l'API Render ont toutefois été vérifiés fonctionnellement par une exécution externe ; les paramètres Vercel/Render et les secrets ne sont pas reproductibles depuis le dépôt seul.
\end{itemize}

Ces éléments peuvent figurer parmi les perspectives (conclusion), sans être présentés comme livrés.

\section{Conclusion}

Ce chapitre a identifié les acteurs du système, traduit la problématique en besoins fonctionnels et non fonctionnels, puis distingué les exigences réalisées des éléments restant hors périmètre. Le chapitre suivant présente la conception de l'architecture, des données et des principaux comportements de la solution.
